\subsection{Which bounds describe a team's work?}

Suppose two engineers each spent one person-hour reviewing the same change
during a two-hour interval. One left after the first hour, so their review
occupied that whole hour. The other engineer could have reviewed at any time
during the two hours. Their combined first-hour labour therefore ranges
from one to two person-hours. If they confirm that they reviewed together
throughout, the first engineer's departure fixes the timing for both. They
contributed exactly two person-hours in the first hour.

The labour bound may also tell a team whether it can report a rate. Suppose
its output score $Q$ is known and Equation~\eqref{eq:sharp-labour-range}
allows any labour total in $[0,H_{\max}]$ with $H_{\max}>0$. For $Q>0$, every rate in
$[Q/H_{\max},\infty)$ is possible among histories with positive labour,
since labour can approach zero. For $Q=0$, the rate is zero in every such history. A history with zero labour has an undefined rate. If
$H_{\max}=0$, report $Q$ alongside zero person-hours. The accompanying
\texttt{bound-labour} command omits the productivity interval whenever
$H_{\min}=0$.

Even with positive labour in every feasible history, the output may
depend on how much labour occurred. Suppose the evidence permits exactly
the pairs $(H,Q)=(h,2h)$ for $1\leq h\leq2$, with $H$ in person-hours
and $Q$ in reference units. At the lowest output and highest labour the ratio would be one unit
per person-hour. Pairing the highest output with the lowest labour would imply four.
In every feasible pair the output is twice the labour,
so the only possible rate is two units per person-hour. More generally,
for a candidate rate $r$ and
positive $H$, the inequality $Q/H\leq r$ is equivalent to $Q-rH\leq0$.
An upper bound of zero on $Q-rH$ over the feasible pairs therefore
establishes $r$ as an upper bound on productivity. If linear constraints
describe the pairs, optimising $Q/H$ over them is a linear-fractional
programme \cite[Section~4.3.2]{boydvandenberghe2004}.
